\documentclass[11pt,a4paper]{article}
\usepackage{amsmath,amssymb,amsthm}
\usepackage[colorlinks=true,linkcolor=blue,citecolor=blue,urlcolor=blue]{hyperref}
\usepackage[margin=1in]{geometry}
\usepackage{xcolor}
\usepackage{booktabs}
\usepackage{array}
\usepackage{float}
\usepackage{mathtools}

\newcommand{\Sp}{\mathrm{Sp}}
\newcommand{\rad}{\mathrm{rad}}
\newcommand{\rank}{\mathrm{rank}}
\newcommand{\F}{\mathbb{F}}
\newcommand{\Z}{\mathbb{Z}}
\newcommand{\Nanti}{N_{\mathrm{anti}}}
\newcommand{\Sq}{\mathrm{Sq}}

% --- drafting flag: visible TODOs for the spec-side anchoring the author must confirm ---
\newcommand{\userflag}[1]{\textcolor{red}{\textsf{[\textbf{USER:} #1]}}}

\theoremstyle{plain}
\newtheorem{theorem}{Theorem}[section]
\newtheorem{lemma}[theorem]{Lemma}
\newtheorem{proposition}[theorem]{Proposition}
\theoremstyle{definition}
\newtheorem{principle}[theorem]{Principle}
\theoremstyle{remark}
\newtheorem*{remark}{Remark}

\title{\textbf{The n/ Program:\\
  Semantic Self-Representation and Its Dimensional Obstruction}\\[4pt]
\large (Paper 0 --- the n/ Whitepaper for CAID)}
\author{Zhou-Li Chen\\[2pt]\small co-nlang Research}
\date{June 2026}

\begin{document}
\maketitle

\begin{abstract}
n/ is a semantic operating system built on a single stance: meaning is
\emph{discovered, not invented}, and a system worth the name must be able to
represent --- to \emph{observe} --- itself coherently. The second half of that
stance is the engineering decision we call \textbf{CAID}: a structure should
carry a coherent self-interpretation. This whitepaper records the mathematical
fate of that decision. Across a twenty-two paper series the same object is
computed at successively higher cohomological degree --- the obstruction to a
structure observing itself --- and the verdict is sharp: \emph{semantic
self-representation is dimensionally obstructed, and it is free of obstruction at
exactly one dimension, $n=4$}. Concretely, the Mermin/anticommutation class
$[\,n_a\,]=[\delta\mu]$ vanishes universally if and only if $n=4$ (Paper~XXI).
The series and the language specification are not two projects: the spec
\emph{poses} the question (can a semantic structure observe itself without
obstruction?) and the papers \emph{answer} it (yes, freely, only at $n=4$). This
document is the bridge between the two sides, and --- like every result it
summarizes --- it is itself a bridge.
\end{abstract}

% =====================================================================
\section{What n/ set out to be (the spec side)}
% =====================================================================

The n/ specification opens with a stance, not a feature list: \emph{writing
programs is discovery, not invention}. The programmer is an \emph{observer} of a
structure that is already there; the work is to find the essence of a thing, not
to legislate it. Three design commitments follow, and they are the engineering
decisions the paper series exists to back:

\begin{itemize}
  \item \textbf{Observation as the ground.} Meaning is accessed only through
  \emph{contexts} --- locally consistent frames of observation. This is the
  language-level form of Bohr's doctrine of classical concepts, and it is the
  reason Bohrification (Paper~II) sits at the foundation rather than as an
  afterthought. \userflag{confirm the spec's precise notion of ``context'' (the
  unit of observation) and how the lattice / smell-search layer presents it.}
  \item \textbf{Essence via relationships (CAID).} A thing is determined by how
  it is observed from everywhere else --- the Yoneda stance. \textbf{CAID} is the
  decision that a structure should carry a \emph{coherent self-interpretation}:
  it should be able to represent itself faithfully inside itself.
  \userflag{fill the exact expansion/definition of CAID from the spec; we use it
  here as ``coherent self-representation = the Yoneda self-embedding.''}
  \item \textbf{Discovery, hence realism --- but about what?} If structure is
  discovered, it is objective. The series will sharpen \emph{what} is objective:
  not observer-independent \emph{values} (those turn out not to exist), but the
  observer-independent \emph{structure of observation} itself.
\end{itemize}

The whitepaper's question is the one the spec poses implicitly the moment it
commits to CAID over a contextual semantics:

\begin{quote}
\emph{Can a semantic structure observe --- represent --- itself coherently, and
if so, when?}
\end{quote}

The obstruction-ladder series is the part of n/'s mathematical backing that
answers this hardest question. The remaining sections give the answer and show
it is forced, not chosen.

% =====================================================================
\section{From semantics to symplectic geometry (the bridge)}
% =====================================================================

This is the load-bearing translation: the dictionary that turns the spec-side
question into a computation. Everything downstream is the evaluation of one
cohomological obstruction along this dictionary.

\begin{center}
\renewcommand{\arraystretch}{1.35}
\begin{tabular}{>{\raggedright}p{0.42\linewidth} p{0.50\linewidth}}
\toprule
\textbf{spec / observation side} & \textbf{formal / symplectic side} \\
\midrule
the field of observables &
$V=\F_2^{2n}$, the phase space over $\F_2$ \\
a \emph{context} (a maximal set of jointly observable quantities) &
a \emph{Lagrangian} $L\subset V$ (maximal isotropic subspace) \\
``can two observables be seen together'' &
commutation $\omega(u,v)=0$ ($\omega$ the symplectic form) \\
the phase / sign carried by an observable &
the quadratic refinement $q$, $q(v)=$ the squaring class \\
the diagram of how contexts overlap &
the \emph{context-nerve} $K_N\simeq S^{N-2}$ \\
CAID --- coherent self-representation &
the Yoneda self-embedding, i.e.\ the realization map\newline
$f:\ \text{nerve}\longrightarrow BH$ \\
the operators with their phases &
the extra-special $2$-group $H$,\newline
$1\to \Z/2\to H\to V\to 1$ \\
``the self-representation is incoherent at degree $k$'' &
a nonzero class in $H^{k}$ (the transgression tower) \\
\bottomrule
\end{tabular}
\end{center}

\noindent Two remarks fix the meaning of the dictionary.

\textbf{Contexts are observation frames.} A Lagrangian is exactly a maximal
collection of observables that can be \emph{seen at once} (pairwise commuting).
The nerve records how these frames overlap; a global section of the nerve would
be an observer-independent assignment of values consistent across all frames.
The whole theory measures the \emph{failure} of such global sections.

\textbf{CAID is the Yoneda self-embedding.} ``A structure represents itself''
means: there is a map from the structure (presented by its nerve of contexts)
into its own classifying object $BH$, recovering the structure from how it is
observed. \userflag{confirm that this is the intended reading of CAID at the spec
level; the math side ($f:\text{nerve}\to BH$) is what the papers actually
compute.} The obstruction to this map being coherent is the content of the
series. Crucially, the primary form of the map is \emph{zero}
($f^*\omega=0$, since $H^2(S^{N-2})=0$ in the relevant range), so coherence is a
\emph{secondary}, higher-coherence datum --- which is precisely why a single
finite formula never closes it, and why the ``bridge'' is rebuilt at every
degree (\S\ref{sec:ladder}, \S\ref{sec:onearena}).

% =====================================================================
\section{The ladder: one bridge at successive degrees}
\label{sec:ladder}
% =====================================================================

Under the dictionary, ``coherent self-representation'' becomes a tower of
cohomological obstructions, each a transgression in the spectral sequence of
$1\to\Z/2\to H\to V\to1$. Kudo's theorem $\tau(\Sq^i x)=\Sq^i\tau(x)$ makes the
bridge \emph{self-replicate} under Steenrod operations: the same construction
reappears at each degree. The rungs, with the paper that establishes each:

\begin{center}
\renewcommand{\arraystretch}{1.3}
\begin{tabular}{p{0.10\linewidth} p{0.40\linewidth} p{0.40\linewidth}}
\toprule
degree & obstruction & paper(s) \\
\midrule
(base) & Bohrification: meaning via contexts (the $1$-truncated bridge)
       & Paper~II \\
$H^2$  & the Maslov / Kashiwara class $\mu$ ($=[\omega]$)
       & Paper~XIII; cf.\ III--V \\
$H^3$  & the anticommutation class $n_a=\Sq^1\omega=\delta\mu$ \emph{(the rung
         that opens)} & Papers~IX, XX \\
$H^3$  & vanishing $\iff n=4$ (the master dichotomy)
       & Papers~XVIII, XXI \\
$H^4$  & the resonance ceiling: no genuine arity-$5$ rung (truncation)
       & Paper~XXII \\
\bottomrule
\end{tabular}
\end{center}
\userflag{check the per-rung paper attributions against the actual series (the
filenames suggest II=bohrification, IX=obstruction\_ladder, XIII=maslov\_index,
XVIII=mermin\_sp8, XX/XXI=h3\_modulus/master, XXII=resonance\_ceiling); adjust as
needed.}

The reading is uniform: each rung is the self-representation bridge at one
cohomological degree. The base case is Bohrification itself --- the original,
$1$-truncated bridge that presents a (noncommutative) structure by the diagram
of its (commutative) observation contexts. n/ computes the \emph{higher} gluing
obstructions that Bohrification left implicit. In this sense the series is the
\emph{cohomological completion of Bohrification}.

% =====================================================================
\section{The master theorem: $n=4$}
% =====================================================================

The ladder's central output is a single dichotomy.

\begin{theorem}[Master dichotomy; Papers~XVIII, XXI]
The obstruction vanishes universally,
\[
   [\,n_a\,]=[\delta\mu]=0 \ \text{for all configurations} \iff n=4 .
\]
For $n=3$ self-representation is possible only partially; for $n=4$ it is free of
obstruction; for $n\ge5$ it is genuinely obstructed (the $H^3$ class is alive,
carried by the modulus $\rad(\omega|_W)$).
\end{theorem}

Read back through the dictionary, this is the answer to the whitepaper's
question:

\begin{principle}[CAID is dimensionally bounded]
A semantic structure can coherently represent itself --- carry CAID without
obstruction --- at \emph{exactly one} dimension, $n=4$. The engineering decision
to demand coherent self-representation is therefore not free: it selects $n=4$.
\end{principle}

\noindent The $n=4$ point is not an accident of low-dimensional coincidence; it
is the unique place where the otherwise-bottomless tower of self-observation
\emph{closes at finite degree} (Papers~XVIII--XXI; the even/odd structure of the
$n\ge5$ modulus, Papers~XIX--XX). \userflag{one sentence here on \emph{why $n=4$}
in spec terms --- e.g.\ which n/ construction lives at $2n=8$ --- if the spec
fixes the dimension for an independent engineering reason, that closes the loop.}

% =====================================================================
\section{The two faces are one}
\label{sec:onearena}
% =====================================================================

The spec side and the paper side were never two projects; they are two views of
one claim. Three identifications make this precise.

\textbf{``Discovery, not invention'' is realism about the
observation-structure, not about values.} Contextuality says there is no
observer-independent assignment of values (no global section of the nerve). So
what is discovered-not-invented cannot be a table of pre-existing values --- it
is the \emph{structure of observation} itself: the nerve, the cohomology, the
bridge. The realism survives exactly because it is relocated from values to
observation-structure. This is the firewall that keeps the philosophy from
floating: every claim about ``essence'' is cashed out as a class in $H^*$.

\textbf{CAID $=$ Yoneda $=$ self-observation $=$ the cohomology tower.} To be is
to be observed (Yoneda); a coherent self-interpretation is the structure
observing itself; the obstruction to that is the transgression tower. These are
the same object named in three languages --- engineering (CAID), category theory
(Yoneda), and cohomology ($H^*$).

\textbf{Observation has no bottom; $n=4$ is where the bottom appears.} Because
the self-representation map is secondary (its primary form is zero), fixing one
global observation frame only spawns the next coherence degree:
$\text{Bohr}\to\mu\,(H^2)\to n_a\,(H^3)\to\cdots$. This bottomless reflexivity is
the honest face of the $\infty$ in $\infty$-Yoneda --- ``building bridges
forever'' is not a failure to arrive but the signature of an $\infty$-categorical
self-embedding, where the bridge \emph{is} the content. The master theorem says
this bottomless tower nonetheless \emph{closes coherently at finite degree} at
one and only one dimension: $n=4$.

% =====================================================================
\section{The honest frontier}
% =====================================================================

A whitepaper should end with what is \emph{not} yet settled, at the right
altitude.

\begin{itemize}
  \item \textbf{The ceiling, and whether it is exotic-proof (item~21).} Paper~XXII
  shows the natural obstruction stops at $H^3$. Whether an \emph{exotic}
  arity-$5$ invariant could furnish a new $H^4$ obstruction for $n\ge5$ is
  \emph{reduced}, not closed: to ``Witt + polarization + a representability
  condition + finite verified computations,'' with the unique degree-$4$ ambient
  invariant being the decomposable $q^2$ (family~B). The single residual is a
  naturality/representability condition --- itself the statement that the
  framework's obstruction is the standard contextuality (value-assignment)
  obstruction.
  \item \textbf{The explicit bridge never finitely closes (item~23).} The
  identification $n_a=\langle\Sq^1\omega,[K_5]\rangle$ holds, but its join is a
  \emph{secondary} operation (Kudo transgression): there is no finite cochain
  formula. This is not incompleteness --- it is the $\infty$-coherence of
  self-representation. The logical ground does bottom out, at the proven
  classical theorems of Kudo and Quillen.
\end{itemize}

Both open ends are of the same kind: they are the perpetual, higher-coherence
\emph{re-building of the one bridge}. That the bridge is never finitely finished
is the mathematical shape of self-representation, not a gap in the program.

% --- references: cite by paper number; replace with \bibitem entries as needed ---
\userflag{add the formal bibliography / Zenodo DOIs (XX:\,10.5281/zenodo.20608302,
XXI:\,10.5281/zenodo.20685669, XXII:\,10.5281/zenodo.20685777; fill the rest) and
a pointer to the n/ specification.}

\end{document}
